% =========================================================
% Section: Darboux's Theorem
% =========================================================
% ---------------------------------------------------------

\begin{theorembox}{Theorem (Darboux's Theorem)}
\begin{theorem}[Darboux's Theorem]
\label{thm:darboux}
Let $I := [a,b] \subseteq \mathbb{R}$ and let $f : I \to \mathbb{R}$ be differentiable on $I$. If $k$ is a real number strictly between $f'(a)$ and $f'(b)$, then there exists $c \in (a,b)$ such that
\[
f'(c) = k.
\]
\hyperref[prf:darboux]{Go to proof.}
\end{theorem}
\end{theorembox}

\begin{remark*}[Standard quantified statement]
\[
\forall k\;\bigl(\min(f'(a),f'(b)) < k < \max(f'(a),f'(b))
\;\Rightarrow\; \exists c \in (a,b)\;(f'(c)=k)\bigr).
\]
\end{remark*}

\begin{remark*}[Predicate reading]
The displayed quantified statement is the predicate-level form of $label: its hypotheses name the input conditions, and its conclusion names the asserted structure or relation.
\[
  f\colon A\to\mathbb{R},\\
  \operatorname{Derivative}(D,f,a,\mathbb{R},\mathbb{R}).
\]
\end{remark*}

\begin{remark*}[Interpretation]
The Intermediate Value Theorem says continuous functions have the intermediate
value property. Darboux's theorem says derivatives also have that property even
when the derivative is not continuous. Thus a derivative cannot have a jump
discontinuity. The proof uses differentiability of $f$ to construct an
auxiliary function $g(x)=f(x)-kx$ and then applies the Extreme Value Theorem
and the Interior Extremum Theorem.
\end{remark*}

\begin{dependencies}
\begin{itemize}
  \item \hyperref[def:derivative-at-a-point]{Definition: Derivative at a Point}
  \item \hyperref[thm:extreme-value-theorem]{Theorem: Extreme Value Theorem}
  \item \hyperref[thm:interior-extremum-theorem]{Theorem: Interior Extremum Theorem}
\end{itemize}
\end{dependencies}
